\documentclass[12pt,a4paper]{article}
\usepackage[utf8]{inputenc}
\usepackage{geometry}
\usepackage{titlesec}
\usepackage{hyperref}
\usepackage{fancyhdr}
\usepackage{graphicx}
\usepackage{listings}
\usepackage{xcolor}

% Define Brand Colors
\definecolor{brandTeal}{HTML}{14b8a6}
\definecolor{brandGray}{HTML}{0f172a} 
\definecolor{brandRose}{HTML}{f43f5e}
\definecolor{codeBg}{HTML}{f1f5f9}

% Geometry setup
\geometry{
    left=25mm,
    right=25mm,
    top=25mm,
    bottom=25mm,
}

% Header and Footer
\pagestyle{fancy}
\fancyhf{}
\lhead{\textcolor{brandTeal}{LandlordPro User Guide}}
\rhead{\today}
\cfoot{\thepage}

% Title Formatting
\titleformat{\section}
{\color{brandTeal}\normalfont\Large\bfseries}
{\thesection}{1em}{}

\titleformat{\subsection}
{\color{brandGray}\normalfont\large\bfseries}
{\thesubsection}{1em}{}

\titleformat{\subsubsection}
{\color{brandGray}\normalfont\normalsize\bfseries}
{\thesubsubsection}{1em}{}

% Code listing style
\lstset{
    basicstyle=\ttfamily\small,
    backgroundcolor=\color{codeBg},
    frame=single,
    rulecolor=\color{brandTeal},
    breaklines=true,
    keywordstyle=\color{brandRose},
    commentstyle=\color{gray},
    stringstyle=\color{brandTeal}
}

% Title Page
\title{\textcolor{brandTeal}{\textbf{User \& Support Guide}}\\for\\LandlordPro}
\author{LandlordPro Team}
\date{\today}

\begin{document}

\maketitle
\thispagestyle{empty}
\newpage

\tableofcontents
\newpage

\section{Introduction}
Welcome to **LandlordPro**! This guide will help you install, configure, and use the LandlordPro property management system. This system is designed for **Administrators** and **Property Managers** to manage properties, tenants, and finances.

\section{Installation \& Setup}

\subsection{Prerequisites}
Before you begin, ensure you have the following installed on your system:
\begin{itemize}
    \item \textbf{Docker Desktop} (Recommended)
    \item \textbf{Git} (for cloning the repository)
\end{itemize}

\subsection{Quick Start (Docker)}
The easiest way to get LandlordPro up and running is using Docker.

\begin{enumerate}
    \item \textbf{Clone the Repository}
    \begin{lstlisting}[language=bash]
git clone <repository-url>
cd landlordpro
    \end{lstlisting}

    \item \textbf{Run the Application}
    \begin{lstlisting}[language=bash]
docker-compose up --build
    \end{lstlisting}

    \item \textbf{Access the App}
    Open your browser and navigate to:
    \begin{itemize}
        \item Frontend: \url{http://localhost}
        \item API: \url{http://localhost:3000}
    \end{itemize}
\end{enumerate}

\section{Getting Started}

\subsection{Registration}
1. Navigate to the login page (\url{http://localhost}).
2. Click on "Register".
3. Fill in your full name, email, and password.
4. Select your role (**Manager**). Note: Admin accounts are typically pre-seeded or created by an existing Admin.

\subsection{Login}
1. Enter your registered email and password.
2. Click "Login".
3. You will be redirected to your specific dashboard based on your role:
   \begin{itemize}
       \item \textbf{Managers} $\rightarrow$ \texttt{/manager/dashboard}
       \item \textbf{Admins} $\rightarrow$ \texttt{/admin/dashboard}
   \end{itemize}

\section{Manager Guide}
As a Manager, your dashboard is your command center for property operations.

\subsection{Dashboard Overview}
Upon logging in, you will see a summary of:
\begin{itemize}
    \item Total Properties managed
    \item Active Leases
    \item Occupancy Rate
    \item Pending Actions (e.g., leases expiring soon)
\end{itemize}

\subsection{Managing Properties \& Units}
The system uses a hierarchical structure: \textbf{Property} $\rightarrow$ \textbf{Floors} $\rightarrow$ \textbf{Locals} (Units).

\begin{enumerate}
    \item Navigate to the **Properties** tab.
    \item Click "Add Property" to define a new building.
    \item **Floors \& Basement**: Enter the number of floors and check "Has Basement" if applicable. The system will \textbf{automatically create} these floors for you.
    \item Select the property to **view** the created floors.
    \item Select a floor to add **Locals** (specific rentable units).
\end{enumerate}

\subsection{Tenant Management}
Since tenants do not log in, you must manage their data.
\begin{enumerate}
    \item Go to the **Tenants** tab.
    \item Click "Add Tenant".
    \item Enter details (Name, ID/TIN, Phone, Email).
    \item Assign them to a specific lease using the **Leases** tab.
\end{enumerate}

\subsection{Financials: Payments \& Expenses}
\subsubsection{Recording Payments}
When a tenant pays rent (Cash, Bank Transfer, etc.):
1. Navigate to **Payments**.
2. Click "Record Payment".
3. Select the Tenant and Lease.
4. Enter the Amount and Date.
5. Save. The system will update the lease balance.

\subsubsection{Tracking Expenses}
1. Navigate to **Expenses**.
2. Log costs such as repairs, cleaning, or utilities.
3. Link the expense to a specific property for accurate reporting.

\section{Administrator Guide}
Admins have a global view of the system components.

\subsection{User Management}
*   **Create/Manage Users**: View all registered managers. You can activate or deactivate accounts.
*   **Role Assignment**: Ensure users have the correct permissions (Admin vs Manager).

\subsection{System Configuration}
*   **Payment Modes**: Configure accepted payment types (e.g., "Cash", "Check", "Wire Transfer") via the Settings page.
*   **Global Reports**: View aggregated financial data across all properties and managers.

\section{Troubleshooting}

\subsection{Common Issues}

\textbf{Issue}: "Connection Refused"
\textbf{Solution}: Ensure Docker containers are running. Run:
\begin{lstlisting}[language=bash]
docker-compose ps
\end{lstlisting}

\textbf{Issue}: Forgot Password
\textbf{Solution}: Password resets are currently handled ensuring security. Please contact an Administrator to reset your credentials. Self-service password updates are only available via the Profile page \textbf{after} logging in.

\subsection{Support}
For technical support, please verify logs:
\begin{lstlisting}[language=bash]
docker-compose logs -f backend
\end{lstlisting}
Or contact the development team.

\end{document}
